\documentclass{article}
\usepackage{amsmath}
\usepackage{mathtools}

\begin{document}
	\section*{Task}
	Consider the Lotka-Volterra model with prey logistic growth:
	\begin{enumerate}
		\item Find the system's equilibrium points.
		\item Calculate the Jacobian matrices of the system at the equilibrium points.
		\item Calculate the eigenvalues of the Jacobian matrices.
		\item Discuss the results.
	\end{enumerate}
	
	\section*{Solution}
	Consider the following system of differential equations: \\
	\begin{align}
		&\begin{cases}
			\begin{aligned}
				\frac{dx}{dt} &= rx\left(1 - \frac{x}{K}\right) - a x y \vspace{0.5cm} \\[1em]
				\frac{dy}{dt} &= b x y - d y
 			\end{aligned}
		\end{cases}
	\end{align}
	where \\
	\begin{itemize}
		\item $x\ge0$: Prey population
		\item $y\ge0$: Predator population
		\item $r>0$: Prey growth rate
		\item $K>0$: Prey environment carrying capacity
		\item $a>0$: Predation rate
		\item $b>0$: Predator reproduction rate
		\item $d>0$: Predator death rate
	\end{itemize}
	\subsection*{1. Finding the equilibrium points}
	\begin{align}
		&\begin{cases}
			\begin{aligned}
				\frac{dx}{dt} &= 0 \\[1em]
				\frac{dy}{dt} &= 0
			\end{aligned}
		\end{cases}
		\Rightarrow
		\begin{cases}
			\begin{aligned}
				&rx \left(1 - \frac{x}{K} \right) - axy = 0 \\[1em]
				&bxy - dy = 0
			\end{aligned}
		\end{cases}
		\Rightarrow
		\begin{cases}
			\begin{aligned} 
				&x \left(r \left(1 - \frac{x}{K} \right) - ay \right) = 0 \\[1em]
				&y \left(bx - d \right) = 0
			\end{aligned}
		\end{cases}
	\end{align}
	Therefore, one of the equilibrium points is $(x=0, y=0)$ $\rightarrow$ total extinction. \\ \\
	Now to analyze the case $(x \neq 0, y = 0)$: \\
	\begin{align}
			\begin{aligned} 
				&r \left(1 - \frac{x}{K} \right) - ay = 0 \quad \xRightarrow{y = 0}  \quad r - \frac{rx}{K} = 0 \Rightarrow \\[1em] 
				&\Rightarrow \frac{rx}{K} = r \quad \Rightarrow \quad x = K
			\end{aligned}
	\end{align}
	Another equilibrium point is $(x=K, y=0)$ $\rightarrow$ predator extinction, prey reaches carrying capacity. \\ \\
	Last case, $(x \neq 0, y \neq 0)$: \\
	\begin{align}
		&\begin{cases}
			\begin{aligned}
				&r \left(1 - \frac{x}{K} \right) - ay = 0 \\[1em]
				&bx - d = 0
			\end{aligned}
		\end{cases} \\
		&bx - d = 0 \quad \Rightarrow \quad x = \frac{d}{b} \\ 
		\intertext{By substituting $x$ in the first equation we obtain:}
		&r \left(1 - \frac{d}{bK} \right) - ay = 0 \quad \Rightarrow \quad y = \frac{r(bK - d)}{abK}
	\end{align}
	The final equilibrium point is $(x = \frac{d}{b}, y = \frac{r(bK - d)}{abK})$ $\rightarrow$ coexistence.
	
	\subsection*{2. Calculating the Jacobian matrix}
	\begin{align}
		\intertext{Let $\frac{dx}{dt} = f(x, y)$ and $\frac{dy}{dt} = g(x, y)$. Then the associated Jacobian is:}
		J = \begin{bmatrix}
			\frac{\partial f}{\partial x} & \frac{\partial f}{\partial y} \\[0.5em]
			\frac{\partial g}{\partial x} & \frac{\partial g}{\partial y}
		\end{bmatrix}
		\intertext{By calculating the partial derivatives we obtain:}
		J = \begin{bmatrix}
			r - \frac{2rx}{K} - ay & -ax \\[0.5em]
			by & bx - d
		\end{bmatrix}
	\end{align}
	
	\subsubsection*{2.1. Evaluating the Jacobian at the equilibrium points}
	\textbf{Case 1: ($x=0, y=0$)}
	\begin{align}
		J = \begin{bmatrix}
			r & 0 \\ 0 & -d
		\end{bmatrix}
	\end{align}
	\textbf{Case 2: ($x=K, y=0$)}
	\begin{align}
		J = \begin{bmatrix}
			-r & -aK \\ 0 & bK - d
		\end{bmatrix}
	\end{align}
	\textbf{Case 3: $(x = \frac{d}{b}, y = \frac{r(bK - d)}{abK})$}
	\begin{align}
		J = \begin{bmatrix}
			-\frac{rd}{bK} & -\frac{ad}{b} \\[0.5em]
			\frac{r(bK-d)}{aK} & 0
		\end{bmatrix}
	\end{align}
	
	\subsection*{3. Calculating eigenvalues}
	General formula:
	\begin{align}
		&\det(J - \lambda I) = 0
		\intertext{For $2 \times 2$ matrices:}
		&\lambda^2 -\lambda \operatorname{Tr}(J) + \det (J) = 0 \Leftrightarrow \\[1em] \Leftrightarrow &\lambda_{1,2} = \frac{\operatorname{Tr}(J) \pm \sqrt{\operatorname{Tr}(J)^2 - 4 \det(J)}}{2}
	\end{align}
	\\[1em]
	\textbf{Case 1}
	\begin{align}
		\left.
		\begin{array}{l}
			\operatorname{Tr}(J) = r - d \\
			\det(J) = -rd
		\end{array}
		\right\}
		&\Rightarrow
		\lambda^2 - \lambda(r-d) -rd = 0 \Rightarrow \\
		&\Rightarrow (\lambda-r)(\lambda+d)=0 \notag \Rightarrow \\[1em]
		&\Rightarrow \begin{cases}
			\lambda_1 = r \\
			\lambda_2 = -d
		\end{cases}
	\end{align}
	\\[1em]
	\textbf{Case 2} 
	\begin{align}
		& \det(J - \lambda I) = \begin{vmatrix}
			-r-\lambda & -aK \\
			0 & bK - d - \lambda
		\end{vmatrix} = (-r-\lambda)(bK-d-\lambda) = 0 \Rightarrow \notag \\[1em]
		&\Rightarrow \begin{cases}
			-r-\lambda = 0 \\
			bK-d-\lambda = 0
		\end{cases}
		\Rightarrow 
		\begin{cases}
			\lambda_1 = -r \\
			\lambda_2 = bK - d
		\end{cases} 
	\end{align}
	\\[1em]
	\textbf{Case 3} 
	\begin{align}
		&\left.
		\begin{array}{l}
			\operatorname{Tr}(J) = -\frac{rd}{bK} \\[0.5em]
			\det(J) = rd \left(1 -\frac{d}{bK} \right)
		\end{array}
		\right\} \Rightarrow \lambda^2 + \frac{rd}{bK}\lambda + rd\left(1-\frac{d}{bK}\right) = 0 \Rightarrow \notag \\[1em]
		&\Rightarrow \lambda_{1,2} = \frac{-\frac{rd}{bK} \pm \sqrt{\left(\frac{rd}{bK}\right)^2 - 4rd\left(1-\frac{d}{bK}\right)}}{2} \Rightarrow \\[1em]
		&\Rightarrow \lambda_{1,2} = \frac{1}{2} \left(-\frac{rd}{bK} \pm \sqrt{\frac{r^2 d^2}{b^2 K^2} + \frac{4rd^2}{bK}-4rd} \right)
	\end{align}
	
	\subsection*{4. Discussing the Results}
	\subsubsection*{Case 1}
	\begin{equation}
		\begin{cases}
			\lambda_1 = r > 0 \\
			\lambda_2 = -d < 0
		\end{cases}
	\end{equation}
	The eigenvalues have opposite signs $\Rightarrow$ point $(0, 0)$ is an \emph{unstable} \textbf{saddle point}. Any (infinitesimally) small positive perturbation of the $x$ value (prey population) will lead to prey exponential growth until reaching carrying capacity $K$.
	
	\subsubsection*{Case 2}
	\begin{equation}
		\begin{cases}
			\lambda_1 = -r < 0 \\
			\lambda_2 = bK-d
		\end{cases}
	\end{equation}
	\begin{itemize}
		\item If $bK-d < 0$ i.e. $d > bK$ then $\lambda_2 < 0$, both eigenvalues are negative and the point is a \textbf{stable node}.
		\item If $bK-d > 0$ i.e. $d < bK$ then $\lambda_2 > 0$, eigenvalues have opposite signs and the point is an unstable \textbf{saddle point}. Any small positive perturbation of $y$ will destabilize the system.
		\item If $bK - d = 0$ then $\lambda_2 = 0$, which leads to a \textbf{transcritical bifurcation}.
	\end{itemize}
		
	\subsubsection*{Case 3}
	\begin{equation}
	\lambda_{1,2} = \frac{1}{2} \left(-\frac{rd}{bK} \pm \sqrt{\frac{r^2 d^2}{b^2 K^2} + \frac{4rd^2}{bK}-4rd} \right)
	\end{equation}
	Rewrite $\lambda_{1,2} = \frac{1}{2} (\alpha \pm \sqrt{\beta})$. Then:
	\begin{itemize}
		\item $\alpha < 0 \quad \forall r, d, b, K >0$. 
		\item If $\beta \ge 0$ then the system reaches a \textbf{stable node}.
		\item If $\beta < 0$ then the eigenvalues are complex conjugates with negative real parts, which characterize the point as a \textbf{stable spiral focus} (also called a \emph{spiral sink}).
	\end{itemize}
\end{document}